\noindent\textbf{Background.}
Health measurement in later life has moved beyond disease counts towards quality of life, health expectancy, frailty, intrinsic capacity and capability. These approaches are complementary rather than interchangeable. The unresolved question addressed by this thesis is whether intrinsic health capacity and the social and structural conditions that shape its realisation can be represented within a common, empirically testable health-capability construct.

\noindent\textbf{Objective.}
To develop an axiomatic Capability Realisation Model of health-related capability in later life, together with the integrated, equity-sensitive policy calculation built upon it, and to evaluate it empirically. The model combines Biomedical Potential, social enablement and structural freedom through a non-compensatory realisation function; the study establishes the measurement properties of the domains first and tests the realisation model's prognostic and comparative performance second, separating construct development from outcome prediction.

\noindent\textbf{Methods.}
The study uses the National Health and Aging Trends Study (NHATS), a longitudinal study of older Medicare beneficiaries in the United States. The proposed domains are operationalised from theoretically selected indicators and evaluated using structural equation modelling. Competing measurement structures are considered rather than assuming that a higher-order factor is valid a priori. If a general factor is supported, its score is evaluated prospectively in separate models of five-year all-cause mortality and first unplanned hospitalisation. Prognostic performance is compared with pre-specified measures of frailty and conventional health status. Missing data, complex survey design, factor-score uncertainty, internal validation and subgroup measurement properties are addressed explicitly. The conceptual formulation $(L\cdot Q)S-B$ is treated as a heuristic rather than as the empirical estimator of the latent score.

\noindent\textbf{Results.}
\emph{[To be completed from the final locked analysis. No numerical result or claim of validation should be inserted before the measurement and prognostic analyses have been completed and independently checked.]}

\noindent\textbf{Contribution.}
The central contribution is the axiomatic Capability Realisation Model and the integrated policy calculation built upon it. It is not another weighted index but a realisation function derived from stated axioms, in which non-compensation, biomedical--context complementarity, an exact capability-gap decomposition, and an equity-sensitive, budget-constrained policy gradient are results rather than assertions, and which is falsifiable through a pre-specified head-to-head comparison with an additive composite and the outcome-independent factor score. The measurement and prognostic work provides the domain inputs and the comparator that test the model, and establishes the construct's measurement independently of its subsequent prognostic performance. The $H_{\text{potential}}$ remains a research construct until external validation, recalibration and prospective evaluation demonstrate transportability to other populations, including people ageing with haemophilia.
